\section{Tema 3}
\subsection*{Resolución de Cuestiones}

\begin{enumerate}[label=\textbf{\arabic*.}]

    \item \textbf{La estrategia ``De la granja a la mesa'' tiene como objetivo reducir el peso de la agricultura ecológica en la superficie agrícola utilizada para garantizar el abastecimiento.} \\
    \textbf{Falsa.} \textit{El objetivo de esta estrategia es precisamente aumentar la agricultura ecológica, de forma que llegue a suponer el 25\% de la Superficie Agrícola Utilizada (SAU) para el año 2030.}

    \item \textbf{El diferencial favorable de productividad del trabajo de España respecto a la media de la UE se explica, en parte, por el elevado peso de países con productividades muy bajas y alta cuota de empleo agrario.} \\
    \textbf{Verdadera.}

    \item \textbf{En el marco actual de la PAC, la Comisión Europea ha centralizado todas las decisiones, eliminando la capacidad de los Estados miembros para adaptar las normas a su realidad local.} \\
    \textbf{Falsa.} \textit{La nueva Política Agraria Común (PAC) hace lo contrario: amplía la flexibilidad y la capacidad de decisión de los Estados miembros (a través de los Planes Estratégicos Nacionales) para adaptar las medidas a las condiciones locales.}

    \item \textbf{Los nuevos ``Eco-Esquemas'' de la PAC son requisitos de cumplimiento obligatorio para todos los agricultores que deseen recibir el pago básico.} \\
    \textbf{Falsa.} \textit{Los eco-esquemas son de carácter voluntario; los agricultores quedan en total libertad de acogerse o no a los mismos, aunque la Unión Europea destina un presupuesto obligatorio para financiarlos.}

    \item \textbf{La nueva ``arquitectura verde'' de la PAC es una medida exclusiva destinada a las explotaciones de menor dimensión económica.} \\
    \textbf{Falsa.} \textit{Se trata de un conjunto de medidas generales que engloba el pago básico, los eco-esquemas y otras prácticas medioambientales, y no es exclusiva para pequeñas explotaciones, sino que aplica al conjunto del sector.}

    \item \textbf{Uno de los objetivos históricos de la PAC ha sido garantizar precios elevados para el consumidor final con el fin de maximizar el beneficio agrario.} \\
    \textbf{Falsa.} \textit{Uno de los objetivos fundacionales e históricos del Tratado de Roma para la PAC era asegurar el aprovisionamiento de los mercados a unos ``precios razonables'' para el consumidor, no precios elevados.}

    \item \textbf{El coeficiente de apertura externa de la agricultura española ha experimentado un crecimiento notable desde 1985, reflejando una mayor integración internacional.} \\
    \textbf{Verdadera.} \textit{De hecho, el coeficiente de apertura externa casi se ha triplicado entre 1985 y 2024, evidenciando una fortísima integración comercial.}

    \item \textbf{Italia basa su modelo de productividad agraria en una disponibilidad de superficie por persona ocupada superior a la media de la Europa Verde.} \\
    \textbf{Falsa.} \textit{Italia basa su modelo en una elevada productividad de la tierra (intensividad), mientras que su disponibilidad de superficie agrícola por persona ocupada (ratio de estructuras) es muy reducida en comparación con la media europea.}

    \item \textbf{La gestión de la PAC se simplificó mediante la creación de una Organización Común de Mercados única para todos los productos agrarios.} \\
    \textbf{Verdadera.} \textit{A partir del año 2007, en un esfuerzo de simplificación administrativa, se fusionaron 21 OCM sectoriales en una única OCM para todos los productos agrarios.}

    \item \textbf{La mecanización y sustitución de trabajo por capital a partir de 1960 fue un proceso impulsado, en gran medida, por el éxodo rural hacia las ciudades.} \\
    \textbf{Verdadera.} \textit{La masiva salida de mano de obra del campo encareció el trabajo agrario, lo que forzó a los titulares de las explotaciones a invertir en maquinaria (capital) para sustituir el factor trabajo.}

    \item \textbf{En la estructura productiva española, las pequeñas explotaciones son las que presentan una mayor especialización en orientaciones intensivas como la horticultura o la avicultura.} \\
    \textbf{Falsa.} \textit{Los datos estructurales demuestran que las explotaciones más grandes (generalmente de más de 60 UDE) son las que alcanzan una mayor Producción Estándar Total por hectárea, indicando una mayor intensidad productiva.}

    \item \textbf{Históricamente, los precios que perciben los agricultores por sus productos han crecido por encima del Índice de Precios de Consumo (IPC).} \\
    \textbf{Falsa.} \textit{Históricamente, se ha dado la tendencia contraria: los precios percibidos en origen por los agricultores han crecido a un ritmo inferior al Índice de Precios de Consumo general.}

    \item \textbf{Los productos cárnicos y el ganado en pie forman parte de los capítulos más deficitarios del comercio exterior agrario español.} \\
    \textbf{Falsa.} \textit{Las carnes y el ganado en pie son productos que gozan de una elevada competitividad internacional y, de hecho, representan el segundo capítulo más importante de la exportación agroalimentaria española, arrojando saldos muy positivos.}

    \item \textbf{El crecimiento de la productividad del trabajo agrario en España se explica fundamentalmente por la mejora de los rendimientos por hectárea, siendo irrelevante la ratio de estructuras.} \\
    \textbf{Falsa.} \textit{Ambos factores han contribuido a la productividad; sin embargo, la mejora en la ratio de estructuras (superficie por trabajador) ha aportado en mucha mayor medida, especialmente a partir de 1985, por lo que es un factor fundamental, no irrelevante.}

    \item \textbf{Durante el período previo a la transformación de 1960, el sector agrario español se caracterizaba por presentar una elevada relación capital-producto.} \\
    \textbf{Falsa.} \textit{Antes de 1960, la agricultura española era tradicional y atrasada, caracterizándose por el uso intensivo de mano de obra barata y una muy baja capitalización (escasa relación capital-producto).}

\end{enumerate}

